\section{Symmetry breaking and specialisation}
\label{sec:symmetry}

The permutation symmetry of a hidden layer is the cleanest collective degree of
freedom available in a neural network. In the symmetric phase every hidden unit
carries a similar overlap with every teacher unit; in the specialised phase each
unit commits to one, and the gap \eqref{eq:specgap} becomes non-zero. Exact
results exist for the two-layer committee machine in the proportional limit.
This block asks what survives when depth is added and when the student width no
longer matches the teacher's.

\subsection{The committee machine: a diverging response on a shifted path}

\manuscriptfigure{figures/family_specialization.pdf}{%
  \textbf{Specialisation of a soft committee machine.}
  (A) the specialisation gap against the sample ratio, with $95\%$ outer-seed
  intervals; (B) the disorder susceptibility \eqref{eq:chi-binder}, whose
  interior maximum is marked per size; (C) the Binder cumulant; (D) the
  attempted finite-size collapse. Colour encodes system size on a single-hue
  ramp, light to dark.%
}{fig:specialization}

\paragraph{Observation.}
Across \AtlasSpecializationConditions{} conditions spanning
$\alpha\in\AtlasSpecializationControlRange{}$ and
\AtlasSpecializationSizes{} sizes, the disorder susceptibility has an
admissible interior maximum at three of the six sizes, with peak height growing
as $N^{\AtlasSpecializationSharpening{}}$ over those three. At the other three
the maximum is a single-point spike --- one seed of eight inflates the variance
at exactly one control value --- and is rejected rather than reported
(\S\ref{sec:gates}). Fitting the drift of the surviving peak locations,
$\alpha_{\rm peak}\propto d^{\,\gamma-1}$, gives
\begin{equation}
  \gamma=\AtlasSpecializationGamma{},\qquad
  95\%\ \text{CI}\;\AtlasSpecializationGammaCI{},
  \label{eq:gamma-spec}
\end{equation}
which \emph{includes} the proportional value $\gamma=1$. The verdict is
\emph{\AtlasSpecializationVerdict{}} at grade \AtlasSpecializationGrade{}.

\paragraph{Interpretation.}
The response is critical in the sense the design was built to detect: over the
sizes where the peak survives the spike test, its height grows with $N$ rather
than saturating. Nothing about the boundary's \emph{position} is established.

This family is worth recounting because it is where the analysis and the
authors disagreed twice. On a first pass, with a coarse control grid and a
fixed epoch budget, it produced a confident transition claim that the
consistency requirement rejected. Repaired --- a grid of
\AtlasSpecializationConditions{} conditions and training to a convergence
criterion --- it produced an apparently strong second result: peaks at all six
sizes and a drift exponent whose interval excluded $\gamma=1$, which we wrote up
as evidence that the specialisation boundary does not lie on the proportional
path $n\propto d$. Adding the spike test then removed half the peaks and widened
the interval to \eqref{eq:gamma-spec}, which no longer excludes anything. We
report the weaker result. The stronger one was an artefact of counting
single-seed outliers as critical peaks, and it is the kind of claim this paper
exists to avoid making.

\paragraph{Not established.}
The location, and now also the scaling path. The drift extrapolation
\eqref{eq:drift} and the collapse disagree by more than the consistency
tolerance of \S\ref{sec:gates}, so the signature count is capped at one. We also
tried the obvious follow-up --- redoing the collapse in the rescaled variable
$n/d^{\gamma}$, with $\gamma$ fixed in advance by the peak locations --- and it
did \emph{not} improve on the collapse in $\alpha$.

\paragraph{Next falsifier.}
More seeds before more sizes. The binding constraint here is that eight seeds
are not enough to estimate a variance cleanly: half the size-curves in this
family carry a single-point artefact. At twelve or twenty-four seeds the spike
test should reject far fewer peaks, and \eqref{eq:gamma-spec} would then be
fitted on six points rather than three --- which is what would settle whether
$\gamma$ differs from one.

\subsection{Depth: the boundary leaves the window}

\paragraph{Observation.}
Replacing the two-layer student with deeper networks at matched width and
sweeping the same $\alpha$ grid gives verdict
\emph{\AtlasDepthLadderVerdict{}} over \AtlasDepthLadderConditions{}
conditions.

\paragraph{Interpretation.}
The order parameter still responds --- the deep students recover the teacher ---
but the susceptibility maxima fall at the edge of the sampled control range for
most depths, so the analysis censors them rather than naming an endpoint as a
critical point. Depth appears to push the boundary outward in $\alpha$, which is
consistent with the shifted scaling path found above, but the present grid
cannot separate that from a simple loss of resolution.

\paragraph{Not established.}
Nothing about how the transition depends on depth. This cell is honestly empty.

\paragraph{Next falsifier.}
Re-run the depth ladder on the rescaled control $n/d^{\gamma}$ using the
exponent from \eqref{eq:gamma-spec}. If the peaks return to the interior of the
window, the depth dependence becomes measurable in one sweep.

\subsection{Multi-index targets and width mismatch}

\paragraph{Observation.}
For a $K$-index teacher with a quadratic link, subspace recovery
\eqref{eq:subspace} is measured against $\alpha$ for a matched student and for
one of twice the width, over \AtlasMultiIndexConditions{} conditions. The
verdict is \emph{\AtlasMultiIndexVerdict{}}.

\paragraph{Interpretation.}
Recovery of the planted subspace improves with sample ratio and the response is
well outside seed noise, but no finite-size signature is admissible: there is no
interior susceptibility peak that sharpens, no consistent Binder crossing, and
no collapse that beats the size-shuffled null. Within this window, subspace
recovery behaves like a crossover rather than a transition.

\paragraph{Not established.}
The sequential ``staircase'' by which multi-index directions are acquired one
after another is not resolved here; our observable is the aggregate subspace
overlap, which is insensitive to the order in which directions are learned.

\paragraph{Next falsifier.}
Measure per-direction overlaps rather than the aggregate, and look for
separate thresholds. If individual directions show interior peaks that the
aggregate washes out, the crossover verdict is an artefact of the observable
rather than a property of the system.
